\setchapterpreamble[u]{\margintoc}
\chapter{T-SQL: Definição de Dados (DDL)}
\labch{cet615-ddl}

Este capítulo cobre a terceira apresentação da primeira
aula~\cite{BaiocoCET615}: tipos de dados, tabelas, restrições de
integridade e índices. É o projeto físico na prática: transformar o
esquema lógico em um script DDL e executá-lo no SQL Server. Os exemplos
usam um banco de biblioteca, criado para estas anotações.

\section{Tipos de dados}

\begin{table}[h]
	\begin{tabular}{lll}
		\toprule
		Categoria & Tipos & Tamanho \\
		\midrule
		Caractere & \texttt{char(n)}, \texttt{varchar(n)} & até $n$ bytes \\
		Numérico exato & \texttt{decimal(p,e)}, \texttt{numeric(p,e)} & depende de $p$ \\
		Numérico aproximado & \texttt{float}, \texttt{real} & 8 e 4 bytes \\
		Inteiro & \texttt{int}, \texttt{smallint}, \texttt{tinyint} & 4, 2 e 1 byte \\
		Monetário & \texttt{money}, \texttt{smallmoney} & 8 e 4 bytes \\
		Data e hora & \texttt{datetime}, \texttt{smalldatetime} & 8 e 4 bytes \\
		Binário & \texttt{binary(n)}, \texttt{varbinary(n)} & $n$ bytes \\
		Outros & \texttt{bit} & 1 bit \\
		\bottomrule
	\end{tabular}
	\caption{Tipos básicos do SQL Server apresentados em aula.}
\end{table}

Em \texttt{numeric(p,e)}, $p$ é o total de dígitos e $e$, quantos ficam
depois da vírgula: \texttt{numeric(6,2)} vai até 9999,99. Diferente de
\texttt{float} e \texttt{real}, os tipos exatos não têm erro de
arredondamento, e por isso servem para valores monetários.

\subsection{\texttt{char} ou \texttt{varchar}?}

Foi o ponto mais longo da aula sobre tipos\sidecite{BaiocoCET615Aulas}.
Quando todas as colunas têm tamanho fixo, o SGBD sabe exatamente onde
começa cada coluna do registro. Basta uma coluna \texttt{varchar} para o
registro passar a ter tamanho variável: o tamanho de cada valor precisa
ser gravado e lido antes do próprio valor, e cada \texttt{varchar} a mais
soma esse custo. A regra da professora é usar \texttt{varchar} só quando
o ganho supera o custo: quando o tamanho varia muito (de 1 a 100
caracteres, por exemplo) ou quando a coluna é opcional e quase sempre
vazia. Nas palavras dela, tudo o que é muito flexível tem um custo; se
\texttt{varchar} fosse sempre melhor, \texttt{char} não existiria.

Ela também lembrou que \texttt{money} e os tipos de data e hora dependem
da configuração de idioma do servidor.

\subsection*{Complemento: tipos obsoletos e tipos novos}

A tabela dos slides inclui tipos que a Microsoft marca como obsoletos:

\begin{itemize}
	\item \texttt{text}, \texttt{ntext} e \texttt{image} serão removidos em
	versão futura; a documentação manda usar \texttt{varchar(max)},
	\texttt{nvarchar(max)} e \texttt{varbinary(max)}\sidecite{MSTextImage};
	\item \texttt{timestamp} não guarda data nem hora: é sinônimo de
	\texttt{rowversion}, um contador binário que muda a cada alteração da
	linha. A sintaxe \texttt{timestamp} também está obsoleta, e para
	registrar data e hora a recomendação é
	\texttt{datetime2}\sidecite{MSRowversion}.
\end{itemize}

Desde o SQL Server 2008 existem também \texttt{date} (só data),
\texttt{time} (só hora), \texttt{datetime2} (maior precisão) e
\texttt{datetimeoffset} (com fuso). O esquema da primeira avaliação usa
\texttt{date}.

\section{Tabelas}

\subsection{\texttt{create table}}

\begin{lstlisting}[caption={Tabela com nulidade explícita e valor padrão.}]
use biblioteca
go

create table leitor (
codleitor int not null,
nome char(40) not null,
email varchar(80) null,       -- opcional e de tamanho muito variável
cidade char(30) not null default 'Campinas',
primary key (codleitor)
)
\end{lstlisting}

A tabela é criada no banco em uso (\texttt{use}). Cada coluna pode ser
\texttt{null} ou \texttt{not null} e pode ter um \texttt{default}.

O padrão de nulidade, quando nada é declarado, depende da configuração
do banco no servidor. Por isso a professora recomenda declarar
\texttt{null} ou \texttt{not null} em \emph{toda} coluna: o script pode
rodar em outro servidor, com outra configuração. Nomes de coluna não
levam espaço.

\subsection{\texttt{alter table}}

\begin{lstlisting}[caption={Alterações de estrutura.}]
alter table leitor alter column nome char(60) not null  -- muda o tipo
alter table leitor add telefone char(11) null             -- nova coluna
alter table leitor add default 'SP' for uf                -- novo default
alter table leitor drop column telefone                   -- remove coluna

exec sp_rename 'leitor', 'leitores'           -- renomeia tabela
exec sp_rename 'leitores.cidade', 'municipio' -- renomeia coluna
\end{lstlisting}

Regras e alertas:

\begin{itemize}
	\item a coluna nova entra no fim da tabela e, segundo a aula, precisa
	aceitar nulo: as linhas que já existem não teriam valor para ela. O SQL
	Server também aceita uma coluna \texttt{not null} nova se ela tiver
	\texttt{default}; as linhas existentes recebem o valor padrão
	(\texttt{alter table livro add qtd\_exemplares int not null default
	0});
	\item uma coluna só sai se nada depender dela (restrição, índice,
	\texttt{unique}, \texttt{default});
	\item alterações seguras são as que não prejudicam dados existentes:
	aumentar o tamanho de um \texttt{char}, acrescentar coluna ou
	\texttt{default};
	\item renomear tabela ou coluna de um banco em uso pode quebrar as
	aplicações que usam o nome antigo. Vale pensar nos nomes antes;
	\item como no \refch{cet615-sqlserver}, \texttt{alter} é manutenção de
	banco em uso. Em desenvolvimento, a professora prefere apagar e
	recriar.
\end{itemize}

\subsection{\texttt{drop table}}

\begin{lstlisting}
drop table emprestimo
drop table emprestimo, exemplar, livro   -- várias de uma vez
\end{lstlisting}

Uma tabela referenciada por chave estrangeira não pode ser apagada antes
da tabela que a referencia. A ordem de remoção é a inversa da ordem de
criação.

\section{Restrições de integridade}

As restrições de entidade, de chave e referencial do
\refch{cet615-conceitos} viram cláusulas do \texttt{create table}.

\begin{description}
	\item[\texttt{primary key}:]\index{primary key@\texttt{primary key}} a
	coluna é sempre \texttt{not null}, e o SQL Server cria um índice
	\emph{clusterizado} para ela (o padrão pode ser trocado para não
	clusterizado).
	\item[\texttt{unique}:] chave secundária. Pode haver várias por tabela,
	a coluna aceita nulo, e o SQL Server cria um índice \emph{não
	clusterizado} único para ela.
	\item[\texttt{foreign key}:] a tabela referenciada precisa existir
	antes. O SQL Server \emph{não} cria índice para chave estrangeira.
\end{description}

\begin{lstlisting}[caption={Esquema de biblioteca com as três restrições.}]
create table autor (
codautor int not null,
nome char(40) not null,
primary key (codautor)
)

create table livro (
codlivro int not null,
titulo char(60) not null,
isbn numeric(13,0) not null,
codautor int not null,
primary key (codlivro),
unique (isbn),
foreign key (codautor) references autor
)

create index ix_livro_autor on livro (codautor)  -- índice da FK

create table emprestimo (
codemprestimo int not null,
codlivro int not null,
codleitor int not null,
data date not null,
primary key (codemprestimo),
foreign key (codlivro) references livro,
foreign key (codleitor) references leitor
)

create index ix_emprestimo_livro on emprestimo (codlivro)
create index ix_emprestimo_leitor on emprestimo (codleitor)
\end{lstlisting}

Pontos destacados em aula:

\begin{itemize}
	\item Sem chave primária, chave secundária e, principalmente, chave
	estrangeira, o banco deixa de ser relacional. Nas palavras dela, para
	criar tabela sem restrição, melhor usar um banco não relacional.
	\item \texttt{references autor}, sem coluna, aponta para a chave
	primária, como manda o modelo, e falha se a primária tiver sido
	esquecida. Indicar a coluna (\texttt{references autor (codautor)})
	permite apontar para qualquer coluna \texttt{unique}.
	\item O índice da chave estrangeira é obrigatório na prática: o SQL
	Server não o cria. A explicação da professora é que o índice da
	chave estrangeira normalmente não é único, já que quase toda chave
	estrangeira vem de um relacionamento 1:N ou N:M. Num 1:1, o lado 1
	precisa de índice único ou \texttt{unique} na chave estrangeira.
	\item A ordem do script importa: primeiro as tabelas referenciadas,
	depois as que referenciam.
	\item A professora declara as restrições no fim da tabela. Há uma
	armadilha em declará-las na linha da coluna quando a chave é composta:
	duas \texttt{primary key} dão erro, mas \texttt{unique} e
	\texttt{foreign key} não dão erro e criam uma restrição separada por
	coluna, em vez de uma restrição composta.
\end{itemize}

\subsection{Nomear, remover e acrescentar restrições}

Dar nome à restrição é opcional; sem nome, o SQL Server gera um, que
\texttt{sp\_help tabela} mostra. O nome é necessário para remover:

\begin{lstlisting}
exec sp_help emprestimo                  -- descobre o nome
alter table emprestimo drop constraint FK__emprestim__codli__1234

alter table leitor
add constraint uq_leitor_email unique (email)  -- nome explícito
\end{lstlisting}

Restrições normalmente nascem no \texttt{create table}. Criar uma
chave primária depois que o banco está no ar não faz sentido, porque
nulos e repetidos já podem ter entrado. O \texttt{alter table ... add}
serve para corrigir uma chave secundária esquecida (correndo o risco de
já haver repetidos) ou para criar as tabelas em qualquer ordem e
declarar todas as restrições no fim. É assim que o SSMS gera o script de
um banco existente.

\subsection*{Complemento: ações referenciais e \texttt{check}}

\textcite[seção 4.2]{Elmasri2011} trata dois recursos que a aula não
cobre:

\begin{description}
	\item[Ações referenciais.] Por padrão, o SGBD rejeita a operação que
	violaria a integridade referencial: não se apaga um autor que tem
	livros. A chave estrangeira pode declarar outra ação, separada para
	\texttt{on delete} e \texttt{on update}: \texttt{cascade} (propaga a
	exclusão ou a nova chave), \texttt{set null} ou \texttt{set default}.
	O livro sugere \texttt{cascade} para tabelas que mapeiam
	relacionamentos, atributos multivalorados e entidades fracas. No SQL
	Server, o padrão se chama \texttt{no action}.
	\item[\texttt{check}.] Restrição sobre os valores de uma linha, como
	\texttt{check (preco > 0)}. É a forma declarativa das regras
	simples; regras que envolvem outras tabelas ficam para os gatilhos
	(\refch{cet615-gatilhos}).
\end{description}

\begin{lstlisting}
create table exemplar (
codexemplar int not null,
codlivro int not null,
estado char(1) not null check (estado in ('N', 'U', 'D')),
primary key (codexemplar),
foreign key (codlivro) references livro on delete cascade
)
\end{lstlisting}

\section{Índices}

Um \textbf{índice}\index{índice} é uma estrutura que acelera o acesso às
linhas a partir dos valores de uma ou mais colunas. Ele guarda uma
cópia ordenada dos valores da coluna com ponteiros para as linhas.
Ajuda consultas e também atualizações e exclusões, que precisam
encontrar a linha antes de mudá-la. Em troca, ocupa disco e precisa ser
atualizado a cada \texttt{insert}, \texttt{update} e \texttt{delete}.
Muitos índices inúteis pioram as atualizações sem ganho nas consultas.

O uso é transparente: ninguém pede ao SGBD que use o índice. O
\textbf{otimizador de consultas} decide, e pode ignorar um índice que
não ajuda.

\subsection{Quando usar}

\begin{itemize}
	\item colunas usadas no \texttt{where} de \texttt{select},
	\texttt{update} e \texttt{delete}, desde que a condição devolva
	\emph{poucas} linhas;
	\item colunas usadas em \texttt{order by};
	\item colunas usadas em junções.
\end{itemize}

E quando não usar: em tabelas pequenas e em colunas com poucos valores
distintos. Uma coluna com três valores possíveis devolve, em média, um
terço das linhas por consulta, e a varredura sequencial sai mais barata.
A regra empírica da professora é que, acima de 20\% a 25\% das linhas, a
varredura compensa; tabela pequena, para ela, tem menos de mil linhas.

Resumo da aula: chave primária, chave secundária e chave estrangeira se
indexam ``sem pensar''. O que sobra são índices não únicos e não
clusterizados em colunas buscadas com frequência e com poucos repetidos,
como nome de cliente ou título de livro. Uma coluna como UF ou CEP
repete demais para valer a pena.

\subsection{Clusterizado e não clusterizado}

\begin{description}
	\item[Clusterizado (\emph{clustered}):]\index{índice clusterizado} as
	folhas da árvore B são as próprias páginas de dados, na ordem do
	índice. Por isso só pode haver um por tabela: os dados não podem estar
	fisicamente ordenados por duas colunas ao mesmo tempo. Convém criá-lo
	antes dos outros, porque sua criação reordena a tabela e reconstrói os
	demais índices. A escolha deve ser a coluna mais usada em
	\texttt{where}, junções, \texttt{order by} e \texttt{group by}; na
	falta de melhor, a chave primária.
	\item[Não clusterizado (\emph{nonclustered}):] as folhas guardam
	ponteiros para as linhas, que seguem a ordem do clusterizado. Pode
	haver vários por tabela, e é o tipo padrão do \texttt{create index}.
\end{description}

\begin{figure*}[h]
\centering
\begin{tikzpicture}[
	pag/.style={draw=bdroxo, fill=bdroxoclaro, minimum width=14mm, minimum height=6mm, font=\scriptsize},
	dado/.style={draw=bdazul, fill=bdazulclaro, minimum width=14mm, minimum height=6mm, font=\scriptsize},
	lig/.style={draw=bdcinza}]
	% clusterizado
	\node[pag] (r) at (2.25,0) {1 | 40 | 80};
	\node[pag] (i1) at (0,-1.1) {1 | 20};
	\node[pag] (i2) at (2.25,-1.1) {40 | 60};
	\node[pag] (i3) at (4.5,-1.1) {80 | 95};
	\foreach \n in {i1,i2,i3} \draw[lig] (r.south) -- (\n.north);
	\node[dado] (d1) at (0,-2.3) {linhas 1--39};
	\node[dado] (d2) at (2.25,-2.3) {linhas 40--79};
	\node[dado] (d3) at (4.5,-2.3) {linhas 80--};
	\foreach \a/\b in {i1/d1,i2/d2,i3/d3} \draw[lig] (\a.south) -- (\b.north);
	\node[rotulo, align=center] at (2.25,-3.1) {clusterizado: as folhas \emph{são}\\as páginas de dados, em ordem};
	% não clusterizado
	\begin{scope}[xshift=8cm]
	\node[pag] (s) at (2.25,0) {Ana | Lia};
	\node[pag] (j1) at (0.9,-1.1) {Ana | Bia};
	\node[pag] (j2) at (3.6,-1.1) {Lia | Rui};
	\draw[lig] (s.south) -- (j1.north);
	\draw[lig] (s.south) -- (j2.north);
	\node[dado] (e1) at (0,-2.3) {pág. 7};
	\node[dado] (e2) at (1.5,-2.3) {pág. 2};
	\node[dado] (e3) at (3,-2.3) {pág. 9};
	\node[dado] (e4) at (4.5,-2.3) {pág. 4};
	\draw[seta] (j1.south) -- (e1.north);
	\draw[seta] (j1.south) -- (e3.north);
	\draw[seta] (j2.south) -- (e2.north);
	\draw[seta] (j2.south) -- (e4.north);
	\node[rotulo, align=center] at (2.25,-3.1) {não clusterizado: as folhas \emph{apontam}\\para linhas em outra ordem};
	\end{scope}
\end{tikzpicture}
\caption{Índice clusterizado (pela chave primária) e não clusterizado
(pelo nome). Só pode haver um clusterizado, porque os dados só têm uma
ordem física.}
\end{figure*}

\subsection{Único e composto}

\begin{description}
	\item[Único:] os valores não se repetem; a criação falha se já houver
	repetidos. Índice único em coluna \texttt{not null} equivale, na
	prática, a uma chave primária.
	\item[Composto:] formado por várias colunas, em ordem. Funciona como
	ordenar pela primeira coluna e, nos empates, pela segunda. Por isso
	(\texttt{uf}, \texttt{cidade}) é diferente de (\texttt{cidade},
	\texttt{uf}), e o otimizador usa o índice mesmo quando só a primeira
	coluna aparece na busca. Se for único, o que não se repete é a
	combinação.
\end{description}

Uma consequência aparece no script da aula: a chave primária composta
de uma tabela de itens (\texttt{num\_pedido}, \texttt{codigo\_p}) já
gera um índice que atende buscas só por \texttt{num\_pedido}. Só a
segunda chave estrangeira precisa de índice próprio. Pelo mesmo motivo,
a professora reconheceu que o exemplo de índice composto do slide, sobre
as mesmas colunas da chave primária, só serve para mostrar a sintaxe.

\subsection{\texttt{create index} e \texttt{drop index}}

\begin{lstlisting}
create [unique] [clustered | nonclustered] index nome
on tabela (coluna1 [, coluna2, ...])

create index ix_livro_titulo on livro (titulo)
create unique index ix_leitor_cpf on leitor (cpf)
create index ix_exemplar_livro_estado on exemplar (codlivro, estado)

drop index livro.ix_livro_titulo
drop index livro.ix_livro_titulo, leitor.ix_leitor_cpf
\end{lstlisting}

O \texttt{drop index} usa \texttt{tabela.indice} porque o mesmo nome de
índice pode existir em tabelas diferentes. Um índice único em
\texttt{cpf} tem o mesmo efeito de declarar \texttt{unique (cpf)} na
tabela.

\subsection*{Complemento: páginas, limites e a teoria dos índices}

A documentação corrige dois números ditos em aula:

\begin{itemize}
	\item o limite de índices não clusterizados por tabela é 999 desde o
	SQL Server 2008; os 249 do slide eram o limite das versões
	anteriores\sidecite{MSCapacity};
	\item a página do SQL Server tem tamanho fixo de 8 KB e não é
	configurável; oito páginas contíguas formam uma \emph{extensão} de
	64 KB, e a leitura e a gravação em disco são feitas por página
	inteira\sidecite{MSPages}. Uma linha comporta até 8.060 bytes. O que
	o DBA ajusta é o \emph{fill factor}, a fração de cada página de índice
	preenchida na criação.
\end{itemize}

Na teoria, \textcite[cap. 18]{Elmasri2011} distingue três índices de um
nível: o \emph{primário}, sobre o campo-chave de um arquivo ordenado por
ele; o de \emph{agrupamento} (\emph{clustering}), sobre um campo não
chave pelo qual o arquivo está ordenado; e o \emph{secundário}, sobre
qualquer outro campo. Um arquivo tem no máximo um índice primário ou de
agrupamento, e quantos secundários quiser. É a mesma restrição do índice
clusterizado do SQL Server. O capítulo também explica as árvores B e
B$^+$, a estrutura dos índices de vários níveis que o SQL Server usa.

\section{Do esquema ao script}

O roteiro da aula para transformar o esquema lógico em DDL:

\begin{enumerate}
	\item criar as tabelas na ordem do esquema lógico, das referenciadas
	para as que referenciam;
	\item em cada tabela, declarar tipos, nulidade, chave primária,
	chaves secundárias e chaves estrangeiras;
	\item criar um índice para cada chave estrangeira que não seja a
	primeira coluna da chave primária;
	\item avaliar índices extras pelas consultas previstas;
	\item conferir o resultado no diagrama do SSMS (\emph{Database
	Diagrams} $\rightarrow$ \emph{New Database Diagram}), que desenha o
	esquema físico a partir do catálogo;
	\item para desfazer, apagar as tabelas na ordem inversa.
\end{enumerate}
